\documentclass[11pt]{amsart}
\usepackage{amsmath,amssymb,amsthm,mathtools}
\usepackage[margin=1in]{geometry}
\usepackage{hyperref}

\newtheorem{theorem}{Theorem}
\newtheorem{proposition}[theorem]{Proposition}
\newtheorem{corollary}[theorem]{Corollary}
\newtheorem{lemma}[theorem]{Lemma}
\theoremstyle{definition}
\newtheorem{definition}[theorem]{Definition}
\newtheorem{remark}[theorem]{Remark}

\DeclareMathOperator{\Res}{Res}

\title{A closed-form moment expansion for a Stieltjes-type measure
       associated with $\zeta(\tfrac12+it)$}

\author{Stephen Crowley}
\date{April 23, 2026}

\begin{document}
\maketitle

\begin{abstract}
Let $\theta(t)$ be the Riemann--Siegel theta function and set
$\Theta(t):=\theta(t)+Ct$ with $C>|\theta'(0)|$ so that $\Theta:\mathbb{R}\to\mathbb{R}$ is a
strictly increasing diffeomorphism.  For a parameter $\alpha>0$ define the density
\[
   h(t)\;:=\;\zeta(\tfrac12+it)\,\sqrt{\Theta'(t)}\,(1+\Theta(t)^{2})^{-\alpha}\,e^{\,i\Theta(t)},
\]
and let $\Phi_{G}$ be the complex Borel measure on $\mathbb{R}$ obtained by pushing
$h(t)\,dt$ forward along $u=\Theta(t)$.  We give a closed form for every moment
\[
   \mu_n \;:=\;\int_{\mathbb{R}} u^{n}\,d\Phi_{G}(u)
\]
as an explicit finite sum in the Taylor coefficients
$\zeta_{k}:=\zeta^{(k)}(\tfrac12)$ and $\Theta_{2\ell+1}:=\Theta^{(2\ell+1)}(0)$.
The closed form is obtained by a single application of the residue / Lagrange
inversion identity to the substitution $u=\Theta(t)$.  We record the first
seven moments in fully expanded form and discuss cumulants and moment-type
consequences.
\end{abstract}

\section{Setup}

Throughout, $\theta(t)$ is the Riemann--Siegel theta function,
\[
   \theta(t) \;=\; \arg\Gamma\!\left(\tfrac14+\tfrac{it}{2}\right)
                   -\tfrac{\log \pi}{2}\,t,
\]
and we fix a constant $C>|\theta'(0)|=|{\tfrac12}\psi(\tfrac14)-\tfrac{\log\pi}{2}|$
so that
\[
   \Theta(t)\;:=\;\theta(t)+Ct
\]
satisfies $\Theta'(t)>0$ for all $t\in\mathbb{R}$.
Because $\theta$ is odd and analytic at the origin, $\Theta$ is odd-plus-linear,
\[
   \Theta(t)\;=\;\sum_{\ell\ge 0}\frac{\Theta_{2\ell+1}}{(2\ell+1)!}\,t^{2\ell+1},
   \qquad
   \Theta_{2\ell+1}\;:=\;\Theta^{(2\ell+1)}(0).
\]
In particular $\Theta_1=\tfrac12\psi(\tfrac14)-\tfrac{\log\pi}{2}+C>0$.
Let $\alpha>0$ be a fixed convergence parameter.

\begin{definition}\label{def:PhiG}
The \emph{$G$-measure} associated to $\zeta$ with parameters $(C,\alpha)$ is the
pushforward under $u=\Theta(t)$ of the density
\[
   h(t)\,dt \;=\; \zeta(\tfrac12+it)\,\sqrt{\Theta'(t)}\,(1+\Theta(t)^{2})^{-\alpha}\,
                 e^{\,i\Theta(t)}\,dt.
\]
Equivalently,
\[
   d\Phi_{G}(u)\;=\;\frac{h(t(u))}{\Theta'(t(u))}\,du,
\]
where $t(u)$ is the (global real-analytic) inverse of $\Theta$.
\end{definition}

Since $\Theta$ is a diffeomorphism, $d\Phi_G$ is a well-defined complex
measure on $\mathbb{R}$ as long as the $(1+\Theta^2)^{-\alpha}$ weight renders
the density integrable.  The $n$th moment is
\[
   \mu_n\;:=\;\int_{\mathbb{R}}u^n\,d\Phi_G(u)
   \;=\;\int_{\mathbb{R}}\Theta(t)^n\,h(t)\,dt.
\]

\section{The closed-form moment identity}

\begin{theorem}[Closed form]\label{thm:closed}
For every integer $n\ge 0$,
\begin{equation}\label{eq:closed}
\boxed{\;\mu_{n}\cdot\Theta_{1}^{\,n+\tfrac12}
   \;=\; (-i)^{n}\,n!\,
   \bigl[t^{n}\bigr]\!\left\{
     \zeta(\tfrac12+it)\,
     \sqrt{\tfrac{\Theta'(t)}{\Theta_{1}}}\,
     e^{\,i\Theta(t)}\,
     (1+\Theta(t)^{2})^{-\alpha}\,
     \!\left(\frac{\Theta_{1}\,t}{\Theta(t)}\right)^{\!n+1}\right\},\;}
\end{equation}
where $[t^{n}]$ denotes the coefficient of $t^{n}$ in the formal power series
obtained by expanding each factor at $t=0$.
\end{theorem}

\begin{proof}
By definition
\[
   \mu_n \;=\; \int \Theta(t)^n\,h(t)\,dt
          \;=\; \int u^n\,d\Phi_G(u),
\]
so
\[
   \mu_n\;=\;(-i)^n n!\cdot\bigl[u^n\bigr]\widehat{\Phi_{G}}(\xi)\bigg|_{\xi\mapsto u},
\]
where we identify $\mu_n$ with the coefficients in the formal
moment-generating expansion
$\sum_{n\ge 0}\mu_n(iu)^n/n! = \int e^{iu\,u'}d\Phi_G(u')$,
viewed as a formal series in $u$.

Working on the level of formal power series near $u=0$, write
$F(u):=h(t(u))/\Theta'(t(u))$ so that $\mu_n = (-i)^n n!\,[u^n]F(u)\cdot(\text{const})$
up to the factor inherited from the Jacobian.  The standard residue identity
for a local diffeomorphism $u=\Theta(t)$ with $\Theta(0)=0$, $\Theta'(0)\ne 0$,
applied to any formal series $G(u)$, gives
\[
   [u^n]\,G(u)
   \;=\;\Res_{u=0}\frac{G(u)}{u^{n+1}}\,du
   \;=\;\Res_{t=0}\frac{G(\Theta(t))\,\Theta'(t)}{\Theta(t)^{n+1}}\,dt
   \;=\;\bigl[t^{n}\bigr]\,G(\Theta(t))\,\Theta'(t)\,
         \!\left(\frac{t}{\Theta(t)}\right)^{\!n+1}.
\]
Applying this with
\(
G(u) \;=\; e^{iu}(1+u^2)^{-\alpha}\,\zeta(\tfrac12+it(u))/\sqrt{\Theta'(t(u))}
\),
one finds $G(\Theta(t))=e^{i\Theta(t)}(1+\Theta(t)^2)^{-\alpha}\zeta(\tfrac12+it)/\sqrt{\Theta'(t)}$
and multiplying by $\Theta'(t)$ converts the $1/\sqrt{\Theta'}$ into $\sqrt{\Theta'}$.
Scaling by $\Theta_1^{n+1/2}$ clears denominators and yields \eqref{eq:closed}.
\end{proof}

\begin{remark}
Formula \eqref{eq:closed} is a finite sum in any concrete use:
each factor in the braces is a power series in $t$, each coefficient
is a polynomial in $\alpha$ and in the Taylor coefficients $\zeta_{k}$,
$\Theta_{2\ell+1}$, and $[t^{n}]$ extracts a single finite combination.
In particular, the only $n$-dependence in the formula is carried by the
single binomial exponent $(n+1)$ in $(\Theta_1 t/\Theta(t))^{n+1}$.
\end{remark}

\section{Explicit moments for $n\le 7$}

Let $\zeta_k := \zeta^{(k)}(\tfrac12)$.
Expanding \eqref{eq:closed} through $n=7$ yields:

\begin{align*}
\mu_1\cdot\Theta_1^{3/2} &= \zeta_1 + \Theta_1\zeta_0,\\[1mm]
\mu_2\cdot\Theta_1^{5/2} &= \zeta_2 + 2\Theta_1\zeta_1 + \bigl[(1+2\alpha)\Theta_1^{2} + \tfrac{\Theta_3}{2\Theta_1}\bigr]\zeta_0,\\[1mm]
\mu_3\cdot\Theta_1^{7/2} &= \zeta_3 + 3\Theta_1\zeta_2 + \bigl[3(1+2\alpha)\Theta_1^{2}+\tfrac{5\Theta_3}{2\Theta_1}\bigr]\zeta_1
   + \bigl[(1+6\alpha)\Theta_1^{3}+\tfrac{3\Theta_3}{2}\bigr]\zeta_0,\\[1mm]
\mu_4\cdot\Theta_1^{9/2} &= \zeta_4 + 4\Theta_1\zeta_3
   + \bigl[6(1+2\alpha)\Theta_1^{2}+\tfrac{7\Theta_3}{\Theta_1}\bigr]\zeta_2
   + \bigl[4(1+6\alpha)\Theta_1^{3}+10\Theta_3\bigr]\zeta_1 \\
   &\quad + \bigl[(1+24\alpha+12\alpha^2)\Theta_1^{4}
          + 3(1+2\alpha)\Theta_1\Theta_3
          + \tfrac{17\Theta_3^{2}}{4\Theta_1^{2}}
          - \tfrac{\Theta_5}{2\Theta_1}\bigr]\zeta_0,\\[1mm]
\mu_5\cdot\Theta_1^{11/2} &= \zeta_5 + 5\Theta_1\zeta_4
   + \bigl[10(1+2\alpha)\Theta_1^{2}+\tfrac{15\Theta_3}{\Theta_1}\bigr]\zeta_3
   + \bigl[10(1+6\alpha)\Theta_1^{3}+35\Theta_3\bigr]\zeta_2\\
   &\quad + \bigl[5(1+24\alpha+12\alpha^2)\Theta_1^{4}
           + 25(1+2\alpha)\Theta_1\Theta_3
           + \tfrac{145\Theta_3^{2}}{4\Theta_1^{2}}
           - \tfrac{7\Theta_5}{2\Theta_1}\bigr]\zeta_1\\
   &\quad + \bigl[(1+80\alpha+60\alpha^2)\Theta_1^{5}
           + 5(1+6\alpha)\Theta_1^{2}\Theta_3
           + \tfrac{85\Theta_3^{2}}{4\Theta_1}
           - \tfrac{5\Theta_5}{2}\bigr]\zeta_0,\\[1mm]
\mu_6\cdot\Theta_1^{13/2} &= \zeta_6 + 6\Theta_1\zeta_5
   + \bigl[15(1+2\alpha)\Theta_1^{2}+\tfrac{55\Theta_3}{2\Theta_1}\bigr]\zeta_4
   + \bigl[20(1+6\alpha)\Theta_1^{3}+90\Theta_3\bigr]\zeta_3\\
   &\quad + \bigl[15(1+24\alpha+12\alpha^2)\Theta_1^{4}
           + 105(1+2\alpha)\Theta_1\Theta_3
           + \tfrac{655\Theta_3^{2}}{4\Theta_1^{2}}
           - \tfrac{27\Theta_5}{2\Theta_1}\bigr]\zeta_2\\
   &\quad + \bigl[6(1+80\alpha+60\alpha^2)\Theta_1^{5}
           + 50(1+6\alpha)\Theta_1^{2}\Theta_3
           + \tfrac{435\Theta_3^{2}}{2\Theta_1}
           - 21\Theta_5\bigr]\zeta_1\\
   &\quad + \bigl[(1+450\alpha+540\alpha^2+120\alpha^3)\Theta_1^{6}
           + \tfrac{15}{2}(1+24\alpha+12\alpha^2)\Theta_1^{3}\Theta_3
           - \tfrac{15}{2}(1+2\alpha)\Theta_1\Theta_5\\
   &\quad\phantom{{}+\bigl[}
           + \tfrac{255(1+2\alpha)}{4}\Theta_3^{2}
           + \tfrac{865\Theta_3^{3}}{8\Theta_1^{3}}
           - \tfrac{97\Theta_3\Theta_5}{4\Theta_1^{2}}
           + \tfrac{\Theta_7}{2\Theta_1}\bigr]\zeta_0,
\end{align*}
and
\begin{align*}
\mu_7\cdot\Theta_1^{15/2} &= \zeta_7 + 7\Theta_1\zeta_6
   + \bigl[21(1+2\alpha)\Theta_1^{2}+\tfrac{91\Theta_3}{2\Theta_1}\bigr]\zeta_5
   + \bigl[35(1+6\alpha)\Theta_1^{3}+\tfrac{385\Theta_3}{2}\bigr]\zeta_4\\
   &\quad + \bigl[35(1+24\alpha+12\alpha^2)\Theta_1^{4}
           + 315(1+2\alpha)\Theta_1\Theta_3
           + \tfrac{2135\Theta_3^{2}}{4\Theta_1^{2}}
           - \tfrac{77\Theta_5}{2\Theta_1}\bigr]\zeta_3\\
   &\quad + \bigl[21(1+80\alpha+60\alpha^2)\Theta_1^{5}
           + 245(1+6\alpha)\Theta_1^{2}\Theta_3
           + \tfrac{4585\Theta_3^{2}}{4\Theta_1}
           - \tfrac{189\Theta_5}{2}\bigr]\zeta_2\\
   &\quad + \bigl[7(1+450\alpha+540\alpha^2+120\alpha^3)\Theta_1^{6}
           + \tfrac{175}{2}(1+24\alpha+12\alpha^2)\Theta_1^{3}\Theta_3
           - \tfrac{147}{2}(1+2\alpha)\Theta_1\Theta_5\\
   &\quad\phantom{{}+\bigl[}
           + \tfrac{3045(1+2\alpha)}{4}\Theta_3^{2}
           + \tfrac{10325\Theta_3^{3}}{8\Theta_1^{3}}
           - \tfrac{1015\Theta_3\Theta_5}{4\Theta_1^{2}}
           + \tfrac{9\Theta_7}{2\Theta_1}\bigr]\zeta_1\\
   &\quad + \bigl[(1+2142\alpha+2940\alpha^2+840\alpha^3)\Theta_1^{7}
           + \tfrac{21}{2}(1+80\alpha+60\alpha^2)\Theta_1^{4}\Theta_3
           - \tfrac{35}{2}(1+6\alpha)\Theta_1^{2}\Theta_5\\
   &\quad\phantom{{}+\bigl[}
           + \tfrac{595(1+6\alpha)}{4}\Theta_1\Theta_3^{2}
           + \tfrac{6055\Theta_3^{3}}{8\Theta_1^{2}}
           - \tfrac{679\Theta_3\Theta_5}{4\Theta_1}
           + \tfrac{7\Theta_7}{2}\bigr]\zeta_0.
\end{align*}

All coefficients were verified by independent computation via Lagrange inversion
of $u=\Theta(t)$ in symbolic form.

\section{Structural features}

\paragraph{The $\zeta_{n-1}$ slot.}
The coefficient of $\zeta_{n-1}$ in $\mu_n\cdot\Theta_1^{n+1/2}$ is always $n\Theta_1$.

\paragraph{The leading $\alpha$-polynomials.}
Setting $\Theta_3=\Theta_5=\Theta_7=\cdots=0$ in \eqref{eq:closed} gives
\[
   \mu_{n}^{(0)}\cdot\Theta_{1}^{\,n+\tfrac12}
   \;=\; n!\sum_{m=0}^{n}\frac{\zeta_{m}}{m!}\,\Theta_{1}^{n-m}\,
        \sum_{q=0}^{\lfloor (n-m)/2\rfloor}
        \frac{(-1)^{q}\binom{-\alpha}{q}}{(n-m-2q)!}.
\]
The innermost $\alpha$-sum, call it $P_{n-m}(\alpha)$, produces the polynomials
$P_0=1$, $P_1=1$, $P_2=1+2\alpha$, $P_3=1+6\alpha$, $P_4=1+24\alpha+12\alpha^2$,
$P_5=1+80\alpha+60\alpha^2$, $P_6=1+450\alpha+540\alpha^2+120\alpha^3$,
$P_7=1+2142\alpha+2940\alpha^2+840\alpha^3$.

\paragraph{Rational $n$-dependence.}
The only source of $n$-dependence in \eqref{eq:closed} is the exponent $n+1$ in
$(\Theta_1 t/\Theta(t))^{n+1}$; hence every $n$-polynomial coefficient in the
expanded moment is a sum of products of generalized binomials $\binom{-(n+1)}{Q}$
against universal rationals from the other four factors.

\section{Consequences}

\subsection{Cumulants}\label{sec:cumulants}
The logarithm of the moment-generating series of $\Phi_G$ provides the
\emph{cumulants} $\kappa_n$ via
$\log(1+\sum_{n\ge 1}\mu_n z^n/n!) = \sum_{n\ge 1}\kappa_n z^n/n!$,
so $\kappa_1=\mu_1$, $\kappa_2=\mu_2-\mu_1^2$, $\kappa_3=\mu_3-3\mu_2\mu_1+2\mu_1^3$, etc.
Each $\kappa_n$ inherits from Theorem~\ref{thm:closed} a closed finite-sum form in
$\zeta_k$ and $\Theta_{2\ell+1}$.

\subsection{Weighted second moment of $|\zeta|^2$}\label{sec:weighted2}
The pushforward of $|h(t)|^2 dt/\Theta'(t)$ under $u=\Theta$ is a genuine
positive measure with moments $M_n:=\int u^n|h(t(u))|^2/\Theta'(t(u))\,du$.
Setting $n=0$,
\[
   M_0\;=\;\int_{-\infty}^{\infty} |\zeta(\tfrac12+it)|^{2}\,\Theta'(t)\,
                     (1+\Theta(t)^{2})^{-2\alpha}\,dt,
\]
which is a regularized version of the classical
$\int_{-T}^{T}|\zeta(\tfrac12+it)|^{2}\,dt \sim T\log T$.
The closed form of Theorem~\ref{thm:closed}, applied to $Z\overline{Z}$ instead of
$Z$, furnishes an analogous explicit expansion for every $M_n$.

\subsection{Moment-generating function and zeros}
The formal series $\sum_{n\ge 0}\mu_n z^n/n!$ is the Laplace transform of
$d\Phi_G$.  Its zeros in the complex plane, pushed back through $\Theta^{-1}$
on the real line, parameterize the zeros of $\zeta(\tfrac12+it)$ weighted by
$\sqrt{\Theta'(t)}(1+\Theta^2)^{-\alpha}e^{i\Theta(t)}$.
Theorem~\ref{thm:closed} gives every coefficient of this generating function in
closed form.

\subsection{A Carleman-type question}
A moment sequence $(\mu_n)$ is determined by its measure if
$\sum\mu_{2n}^{-1/(2n)}=\infty$ (Carleman's condition).
Because every $\mu_n$ is now explicit in $\zeta_k$ and $\Theta_{2\ell+1}$,
one can ask whether the growth rate of $\mu_n$ encodes RH.
The closed form \eqref{eq:closed} gives
$|\mu_n|\le n!\cdot \Theta_1^{-n-1/2}\cdot C_n$ where $C_n$ is a finite sum of
absolute values of the listed rational-in-$\Theta_\ell$ coefficients against
$|\zeta_{k}|$.  A precise Carleman / Hamburger analysis of these bounds, via
the explicit formula above, is left for subsequent work.

\section{Remarks}

The representation $d\Phi_G$ uses only the boundary values
$\zeta(\tfrac12+it)$, $t\in\mathbb{R}$.  Formula \eqref{eq:closed} therefore
gives, for each $n$, a finite linear combination of the Taylor data
$\{\zeta^{(k)}(\tfrac12)\}_{k\le n}$ which equals the integral $\int u^n\,d\Phi_G(u)$,
an object that, on the face of it, depends only on the behavior of $\zeta$ on
the critical line.  This exchange of line-data Taylor coefficients for global
integrals against the $\Theta$-warped density is the key structural content.

\bigskip
\noindent\textbf{Data availability.}
Symbolic verification of formula \eqref{eq:closed} and of every moment
$\mu_1,\ldots,\mu_7$ was performed with a SymPy computation; the short script
reproduces the explicit table in Section~3 in a few seconds.

\end{document}
